\chapter{Suspension and Expulsion Trial}
\label{sus-exp-trial}

\section{Preface}
	\begin{enumeration}
		\item This section's purpose is to summarize the Suspension and Expulsion Procedures in the \gls{ibl}. For more detailed procedures, refer to \gls{ibl} Article VI , Section 13.
	\end{enumeration}

\section{Purpose}

	\begin{enumeration}
		\item To determine whether suspension or expulsion is warranted due to violations of their obligation to the Fraternity. These violations include:
			\begin{enumeration}
				\item Conduct that indicates he is unworthy to represent, bear the name of, or wear the badge of the Fraternity.
				\item A violation of any oaths taken during Ritual.
				\item A violation of the limitations on memberships as detailed in \gls{ibl} Article VI, Section 12.
				\item Participation in an initiation conducted in volation of \gls{ibl} Article VI, Section 4, Subsection e.
				\item A violation of any financial obligation to the active chapter.
			\end{enumeration}
		\item Suspension should be imposed in cases where an action that provides a path back to good standing is appropriate.
		\item Explusion should be imposed when misconduct warrants permanently extinguishing the member's association with the Fraternity (e.g due to severity, repetition, or other appropriate grounds).
	\end{enumeration}

\section{Accusation}

	\begin{enumeration}
		\item Any collegiate or alumnus member of a chapter may submit an accusation against any other collegiate or alumnus member of that chapter by making a statement of the charges, orally or in writing , privately to the president of the chapter. Each accusation must be as detailed as possible so that the accused may prepare a defense if they so choose.
		\item If the accusing member request action upon the accustation, the president will follow this procedure:
			\begin{enumeration}
				\item Establish a date, time, and place for the trial, which shall be scheduled as soon as reasonably practicable.
				\item Notify the accused in writing of the date, time, and place of the trial, at least seven days prior to the time the trial is scheduled. The notice must contain the specific charges. If personal service is not possible, the president shall notify the accused in writing, by mail, or similar method resulting in delivery of a receipt to the chapter, at least fourteen days prior to the meeting.
				\item Notify all collegiate members of the active chapter of the date, time, and placeof the trial by the fastest means possible.
			\end{enumeration}

	\end{enumeration}

\section{Procedure}

	\begin{enumeration}
		\item After the meeting is called to order and a quorum is determined to be present, the president shall read the accusations, withholding the name of the accuser. The meeting may proceed whether or not the accused member attends.
		\item Witnesses (members or non-members} may be called by any member, including the accused, to testify as to their knowledge of the accusations. Members, including the accused, may ask questions of any witness.
		\item The accused has the right to speak on their own behalf. Members may ask questions of the accused.
		\item After all witnesses have been heard, the accused shall withdraw from the meeting. Members then ballot upon the motion to suspend and/or expel the accused. The affirmative vote of a \gls{simp-maj} of all collegiate members eligible to vote shall be necessary to sustain a motion to suspend. A $\frac{3}{4}$ vote of all collegate members eligible to vote shall be necessary to expel.
		\item If two or more members have been accused of the same or related conduct and the hearings of the accusations for such members have been scheduled for a single meeting, all accused members shall withdraw from the meeting before members vote on any of the otions. No such accused member shall be eligible to vote with respect to a motion relating to any other of the accused.
		\item The result of the votes on all resolutions shall be promptly communicated to the acused member, who shall govern themself accordingly.
		\item No record of the proceeding with respect to an accused member shall be retained unless the motion to suspend or expel is sustained.
			\begin{enumeration}
				\item If the motion to expel is sustained, the chapter shall promptly send \gls{ihq} a copy of the motice of meeting served; a complete statement of the original accusations and any written response made by the accused; a numerical record of the number of al members eligible to vote, the number of members eligible and prsent for the vote, and the number of such members voting for the motion; a summary of all evidence submitted by all witnesses; a copy of the minutes of the meeting recording the action of the chapter.
				\item If a motion to suspend is sustained, the chapter shall promptly report the suspension to \gls{ihq}.
			\end{enumeration}

	\end{enumeration}

\section{Lifting Suspension}
	
	\begin{enumeration}
		\item If the chapter deems that lifting a suspension given by this section is appropriate, the president of the active chapter which imposed the suspension shall receive a petition signed by more than $\frac{1}{3}$ of all collegiate members of the chapter seeking to lift the suspension.
		\item The president shall call and conduct a meeting in accordance with with \gls{ibl} Article VI, Section 13, paragraphs 1 through 6 and the first sentence of paragraph 7.
		\item The members shall ballot upon a motion to lift the suspension. This motion needs the affirmative vote of at least $\frac{1}{2}$ of all collegiate members of the chapter eligible to vote in order to sustain the motion.
	\end{enumeration}

